%!TeX root=Memory.tex

%---------------------------------------------------------
% Resum trilingüe del projecte Demochain
% Català, castellà i anglès
%---------------------------------------------------------


\chapter*{Resum}
\addcontentsline{toc}{chapter}{Resum} \markboth{Resum}{Resum}

Aquest projecte presenta \emph{Demochain}, un sistema de votació descentralitzat de propòsit general construït sobre la
xarxa Algorand i ancorat periòdicament a la xarxa pública Ethereum.
El sistema permet crear organitzacions amb un cens d'adreces autoritzades, definir propostes amb una fase prèvia
d'aprovació per quòrum de dos terços i celebrar eleccions amb vot preferencial.
El recompte final es calcula mitjançant el mètode de Schulze, que garanteix propietats formals com la monotonicitat, la
independència de clons i la consistència amb Condorcet.
Per garantir la integritat del resultat sense dependre d'una única xarxa, un servei d'ancoratge llegeix l'estat final de
la cadena Algorand, calcula un \textit{hash} determinista i el publica conjuntament amb la resta de nodes participants en
el contracte notari \texttt{NotaryContract} desplegat a la \emph{testnet} Sepolia d'Ethereum, on un mecanisme de consens
\textit{K-of-N} fixa el \textit{hash} de manera immutable.
El client web, construït amb React i TypeScript, ofereix verificació pública del vot individual per a qualsevol
observador.
El projecte s'ha desenvolupat amb pràctiques modernes d'enginyeria del \textit{software}: \emph{backlog} jerarquitzat a
GitHub Projects amb èpiques, històries d'usuari i tasques; control de versions amb \emph{feature branches} i revisió
creuada obligatòria; arquitectura documentada amb el model C4; i decisions de disseny recollides en
\emph{Architecture Decision Records}.
L'abast actual constitueix un \ac{MVP} funcional que demostra la viabilitat arquitectònica del plantejament, deixant la
incorporació de privacitat criptogràfica avançada (xifratge homomòrfic ElGamal, proves zk-SNARK i cerimònia de desxifrat
amb \emph{trustees}) com a treball futur.

\textbf{Paraules clau:} \emph{blockchain}, Algorand, Ethereum, votació electrònica, vot preferencial, mètode de Schulze,
ancoratge, \emph{smart contract}, verificabilitat pública.

%---------------------------------------------------------
\vspace{1cm}
\begin{center}
    \textbf{\Large Resumen}
\end{center}
\vspace{0.5cm}

Este proyecto presenta \emph{Demochain}, un sistema de votación descentralizado de propósito general construido sobre la
red Algorand y anclado periódicamente a la red pública Ethereum.
El sistema permite crear organizaciones con un censo de direcciones autorizadas, definir propuestas con una fase previa
de aprobación por quórum de dos tercios y celebrar elecciones con voto preferencial.
El recuento final se calcula mediante el método de Schulze, que garantiza propiedades formales como la monotonía, la
independencia de clones y la consistencia con Condorcet.
Para garantizar la integridad del resultado sin depender de una única red, un servicio de anclaje lee el estado final de
la cadena Algorand, calcula un \textit{hash} determinista y lo publica junto con el resto de nodos participantes en el
contrato notario\texttt{NotaryContract} desplegado en la \emph{testnet} Sepolia de Ethereum, donde un mecanismo de
consenso \textit{K-of-N} fija el \textit{hash} de forma inmutable.
El cliente web, construido con React y TypeScript, ofrece verificación pública del voto individual para cualquier
observador.
El proyecto se ha desarrollado aplicando prácticas modernas de ingeniería del software: \emph{backlog} jerarquizado en
GitHub Projects con épicas, historias de usuario y tareas; control de versiones con\emph{feature branches} y revisión
cruzada obligatoria; arquitectura documentada con el modelo C4; y decisiones de diseño recogidas en \emph{Architecture
Decision Records}.
El alcance actual constituye un \ac{MVP} funcional que demuestra la viabilidad arquitectónica del planteamiento, dejando
la incorporación de privacidad criptográfica avanzada (cifrado homomórfico ElGamal, pruebas zk-SNARK y ceremonia de
descifrado con \emph{trustees}) como trabajo futuro.

\textbf{Palabras clave:} \emph{blockchain}, Algorand, Ethereum, voto electrónico, voto preferencial, método de Schulze,
anclaje, \emph{smart contract}, verificabilidad pública.

%---------------------------------------------------------
\vspace{1cm}
\begin{center}
    \textbf{\Large Abstract}
\end{center}
\vspace{0.5cm}

This project presents \emph{Demochain}, a general-purpose decentralised voting system built on the Algorand blockchain
and periodically anchored to the public Ethereum network.
The system allows users to create organisations with a census of authorised addresses, define proposals that go through
a two-thirds quorum approval phase, and run elections with preferential voting.
The final tally is computed using the Schulze method, which guarantees formal properties such as monotonicity, clone
independence, and Condorcet consistency.
To ensure the integrity of the result without relying on a single network, an anchoring service reads the final state of
the Algorand chain, computes a deterministic \textit{hash}, and submits it together with the other participating nodes
to the \texttt{NotaryContract} smart contract deployed on Ethereum's Sepolia testnet, where a K-of-N consensus mechanism
fixes it immutably.
The web client, built with React and TypeScript, offers public verification of individual votes for any observer.
The project has been developed applying modern software engineering practices: a hierarchical backlog on GitHub Projects
with epics, user stories, and tasks; version control with feature branches and mandatory cross-review; architecture
documented with the C4 model; and design decisions recorded in Architecture Decision Records.
The current scope constitutes a functional \ac{MVP} that demonstrates the architectural viability of the approach,
leaving the integration of advanced cryptographic privacy (ElGamal homomorphic encryption, zk-SNARK proofs, and
threshold decryption ceremony with trustees) as future work.

\textbf{Keywords:} blockchain, Algorand, Ethereum, electronic voting, preferential voting, Schulze method, anchoring,
smart contract, public verifiability.
